% GENERATED FILE. Edit canonical lessons, then run npm run generate:books.
% canonical-source-hash: fnv1a64:4b177549792cc5b3
% canonical-lessons: FR-C193-drap, FR-C193-tiroir, FR-C193-fauteuil, FR-C193-casserole, FR-C193-poele

\chapter{Sheet, Drawer, Armchair, Saucepan, Frying pan}
\label{ch:fr-sheet-drawer-armchair-saucepan-frying-pan}
\hlchaptermodality{193}

\begin{chapteropening}
\textbf{By the end of this chapter:} \emph{I can say a sheet, a drawer, an armchair, a saucepan and a frying pan.}

Complete the last lesson of chapter 193: I can say a sheet, a drawer, an armchair, a saucepan and a frying pan.
\end{chapteropening}

\section[le drap]{le drap — a sheet}

\label{lesson:FR-C193-drap}



\begin{quote}
Before the new one: say the French for clear, light (of a colour), then the French for distant.
\end{quote}

\subsection*{You'll want to know: le drap}
\textbf{le drap} — \textquotedblleft{}a sheet\textquotedblright{}.

\begin{grammarlens}[title={What you've built}]
One more word for this chapter.
\end{grammarlens}

\subsection*{Your turn}
\begin{itemize}
\raggedright
  \item \emph{Say it:} \emph{le drap}
  \item \emph{Say it:} \emph{le drap}, once more
  \item \emph{Say it:} say \emph{le drap}
  \item \emph{Recall:} say the French for clear, light (of a colour), then the French for distant, then say \emph{le drap} again
\end{itemize}

\begin{tcolorbox}[breakable,colback=teal!4,colframe=teal!35!black,title={Before you move on}]
What does le drap mean? (\textquotedblleft{}A sheet\textquotedblright{}.) Say it once more.
\end{tcolorbox}
\section[le tiroir]{le tiroir — a drawer}

\label{lesson:FR-C193-tiroir}



\begin{quote}
Before the new one: say the French for distant, then the French for a sheet.
\end{quote}

\subsection*{You'll want to know: le tiroir}
\textbf{le tiroir} — \textquotedblleft{}a drawer\textquotedblright{}.

\begin{grammarlens}[title={What you've built}]
One more word for this chapter.
\end{grammarlens}

\subsection*{Your turn}
\begin{itemize}
\raggedright
  \item \emph{Say it:} \emph{le tiroir}
  \item \emph{Say it:} \emph{le tiroir}, once more
  \item \emph{Say it:} say \emph{le tiroir}
  \item \emph{Recall:} say the French for distant, then the French for a sheet, then say \emph{le tiroir} again
\end{itemize}

\begin{tcolorbox}[breakable,colback=teal!4,colframe=teal!35!black,title={Before you move on}]
What does le tiroir mean? (\textquotedblleft{}A drawer\textquotedblright{}.) Say it once more.
\end{tcolorbox}
\section[le fauteuil]{le fauteuil — an armchair}

\label{lesson:FR-C193-fauteuil}



\begin{quote}
Before the new one: say the French for a sheet, then the French for a drawer.
\end{quote}

\subsection*{You'll want to know: le fauteuil}
\textbf{le fauteuil} — \textquotedblleft{}an armchair\textquotedblright{}.

\begin{grammarlens}[title={What you've built}]
One more word for this chapter.
\end{grammarlens}

\subsection*{Your turn}
\begin{itemize}
\raggedright
  \item \emph{Say it:} \emph{le fauteuil}
  \item \emph{Say it:} \emph{le fauteuil}, once more
  \item \emph{Say it:} say \emph{le fauteuil}
  \item \emph{Recall:} say the French for a sheet, then the French for a drawer, then say \emph{le fauteuil} again
\end{itemize}

\begin{tcolorbox}[breakable,colback=teal!4,colframe=teal!35!black,title={Before you move on}]
What does le fauteuil mean? (\textquotedblleft{}An armchair\textquotedblright{}.) Say it once more.
\end{tcolorbox}
\section[la casserole]{la casserole — a saucepan}

\label{lesson:FR-C193-casserole}



\begin{quote}
Before the new one: say the French for a drawer, then the French for an armchair.
\end{quote}

\subsection*{You'll want to know: la casserole}
\textbf{la casserole} — \textquotedblleft{}a saucepan\textquotedblright{}.

\begin{grammarlens}[title={What you've built}]
One more word for this chapter.
\end{grammarlens}

\subsection*{Your turn}
\begin{itemize}
\raggedright
  \item \emph{Say it:} \emph{la casserole}
  \item \emph{Say it:} \emph{la casserole}, once more
  \item \emph{Say it:} say \emph{la casserole}
  \item \emph{Recall:} say the French for a drawer, then the French for an armchair, then say \emph{la casserole} again
\end{itemize}

\begin{tcolorbox}[breakable,colback=teal!4,colframe=teal!35!black,title={Before you move on}]
What does la casserole mean? (\textquotedblleft{}A saucepan\textquotedblright{}.) Say it once more.
\end{tcolorbox}
\section[la poêle]{la poêle — a frying pan}

\label{lesson:FR-C193-poele}



\begin{quote}
Before the new one: say the French for an armchair, then the French for a saucepan.
\end{quote}

\subsection*{You'll want to know: la poêle}
\textbf{la poêle} — \textquotedblleft{}a frying pan\textquotedblright{}.

\begin{grammarlens}[title={What you've built}]
One more word for this chapter.
\end{grammarlens}

\subsection*{Your turn}
\begin{itemize}
\raggedright
  \item \emph{Say it:} \emph{la poêle}
  \item \emph{Say it:} \emph{la poêle}, once more
  \item \emph{Say it:} say \emph{la poêle}
  \item \emph{Recall:} say the French for an armchair, then the French for a saucepan, then say \emph{la poêle} again
\end{itemize}

\begin{tcolorbox}[breakable,colback=teal!4,colframe=teal!35!black,title={Before you move on}]
What does la poêle mean? (\textquotedblleft{}A frying pan\textquotedblright{}.) Say it once more.
\end{tcolorbox}
